% Slides — Capítulo 21: Processos Climáticos Naturais
\documentclass[aspectratio=169,11pt]{beamer}
\usepackage{../comum/temameteo}
\graphicspath{{../figuras/cap21/}}

\title[Cap. 21 --- Clima Natural]{Capítulo 21 --- Processos Climáticos
Naturais}
\subtitle{Meteorologia Prática}
\author{}
\date{}

\begin{document}

\begin{frame}
  \titlepage
  \vfill
  \atribuicao
\end{frame}

\begin{frame}{Roteiro}
  \tableofcontents
\end{frame}

% ---------------------------------------------------------------
\section{Feedbacks}

\begin{frame}{O clima como amplificador realimentado}
  \begin{columns}
    \column{0.4\textwidth}
    \eqdestaque{$\Delta T = \dfrac{\lambda_0\,\Delta F}{1 - f}$}
    \vspace{0.4em}
    \begin{itemize}
      \item O diagrama de blocos da eletrônica, com o planeta no
            papel do circuito (T1).
      \item $\lambda_0 = 1/4\varepsilon\sigma T^3$: a resposta de
            Planck, $\sim$1,2\,K por 2$\times$CO$_2$ (T2).
      \item Vapor ($+$, Clausius--Clapeyron!), gelo ($+$), nuvens
            (incerto).
    \end{itemize}
    \column{0.6\textwidth}
    \figlivroslide{fig_21_18e.jpg}{0.46\textheight}{Fig.~21.18 --- o
      laço de realimentação do sistema climático}
  \end{columns}
\end{frame}

\begin{frame}{Gelo-albedo: dois climas para o mesmo Sol}
  \begin{columns}
    \column{0.4\textwidth}
    \begin{itemize}
      \item Frio $\to$ gelo $\to$ albedo $\to$ mais frio: feedback
            positivo $\Rightarrow$ múltiplos equilíbrios (Caso~1).
      \item Ramo quente (nós) e ramo bola-de-neve, com histerese.
      \item O Neoproterozoico visitou o outro ramo — e saiu pelo
            CO$_2$ vulcânico (N1).
    \end{itemize}
    \column{0.6\textwidth}
    \figlivroslide{fig_21_21.jpg}{0.52\textheight}{Fig.~21.21 --- a
      curva em S: limites frio e quente vs.\ insolação}
  \end{columns}
\end{frame}

% ---------------------------------------------------------------
\section{Forçantes naturais}

\begin{frame}{Milankovitch: o metrônomo orbital}
  \begin{columns}
    \column{0.34\textwidth}
    \begin{itemize}
      \item Excentricidade ($\sim$100\,ka), obliquidade (41\,ka),
            precessão (23/19\,ka).
      \item O gatilho: insolação de verão em 65°N (a neve sobrevive
            ao verão?) — T3; Caso~2.
      \item O gelo lento retifica o metrônomo: a serra de 100\,ka
            (N2).
    \end{itemize}
    \column{0.33\textwidth}
    \figlivroslide{fig_21_10.jpg}{0.42\textheight}{Fig.~21.10 ---
      precessão e estações na órbita}
    \column{0.33\textwidth}
    \figlivroslide{fig_21_3.jpg}{0.42\textheight}{Fig.~21.3 --- o
      registro paleoclimático}
  \end{columns}
\end{frame}

\begin{frame}{Vulcões e Sol}
  \begin{columns}
    \column{0.4\textwidth}
    \begin{itemize}
      \item Erupções explosivas põem sulfato na estratosfera:
            $-0{,}5$\,K por 1--2 anos (Pinatubo) — os aerossóis do
            Cap.~7 em escala planetária.
      \item Ciclo solar de 11 anos: $\pm$0,1\% de $S_0$ — pequeno
            comparado ao resto.
      \item São as réguas naturais da atribuição (Cap.~22).
    \end{itemize}
    \column{0.6\textwidth}
    \figlivroslide{fig_21_17b.jpg}{0.48\textheight}{Fig.~21.17 --- a
      transmissão solar despencando após cada grande erupção}
  \end{columns}
\end{frame}

% ---------------------------------------------------------------
\section{Daisyworld}

\begin{frame}{O planeta que se regula sem querer}
  \begin{columns}
    \column{0.34\textwidth}
    \begin{itemize}
      \item Margaridas pretas prosperam com Sol fraco (aquecem);
            brancas com Sol forte (esfriam).
      \item Sem intenção: seleção $+$ acoplamento $=$ feedback
            NEGATIVO emergente (T4; Caso~3).
      \item Diversidade amplia o platô regulado (N3).
    \end{itemize}
    \column{0.33\textwidth}
    \figlivroslide{fig_21_22b.jpg}{0.42\textheight}{Fig.~21.22 --- o
      planeta das margaridas}
    \column{0.33\textwidth}
    \figlivroslide{fig_21_24b.jpg}{0.42\textheight}{Fig.~21.24 ---
      cobertura e temperatura vs.\ luminosidade}
  \end{columns}
\end{frame}

% ---------------------------------------------------------------
\section{O clima presente}

\begin{frame}{Köppen: a circulação geral carimbada no chão}
  \begin{columns}
    \column{0.4\textwidth}
    \begin{itemize}
      \item A tropical (ZCIT) $\cdot$ B seco (subsidência de Hadley)
            $\cdot$ C temperado (frentes) $\cdot$ D continental
            $\cdot$ E polar.
      \item O mapa climático É o mapa dos cinturões do Cap.~11.
      \item Hierarquia de modelos: EBM 0-D $\to$ EBM meridional
            (\texttt{climlab}) $\to$ GCMs (o Cap.~20 rodado por
            séculos).
    \end{itemize}
    \column{0.6\textwidth}
    \figlivroslide{fig_21_25.jpg}{0.52\textheight}{Fig.~21.25 ---
      classificação de Köppen: as letras seguem a circulação}
  \end{columns}
\end{frame}

% ---------------------------------------------------------------
\section{ENSO e a variabilidade interna}

\begin{frame}{Walker, alísios e o gangorra do Pacífico}
  \begin{columns}
    \column{0.3\textwidth}
    \begin{itemize}
      \item Normal: alísios empilham água quente no oeste;
            ressurgência fria no leste (Cap.~11).
      \item El Niño: alísios fraquejam, a piscina quente desliza para
            leste, a convecção vai junto.
      \item Bjerknes: feedback positivo acoplado oceano--atmosfera.
    \end{itemize}
    \column{0.35\textwidth}
    \figlivroslide{fig_21_28a.jpg}{0.42\textheight}{Fig.~21.28 --- La
      Niña: a Walker turbinada}
    \column{0.35\textwidth}
    \figlivroslide{fig_21_28c.jpg}{0.42\textheight}{Fig.~21.28 --- El
      Niño: a convecção migra para leste}
  \end{columns}
\end{frame}

\begin{frame}{O oscilador de recarga}
  \begin{columns}
    \column{0.4\textwidth}
    \eqdestaque{$\dot T = aT + bh$,\quad $\dot h = -cT - dh$
      \ $\Rightarrow$\ 3--5 anos}
    \vspace{0.4em}
    \begin{itemize}
      \item Carga e descarga do conteúdo de calor: oscilação
            irregular excitada pelo ruído do tempo (T5; Caso~4; N4).
      \item Brasil: seca no N/NE e chuva no S em anos de El Niño.
      \item Previsibilidade CLIMÁTICA de meses apesar do caos
            (Cap.~20): o oceano lento é a memória.
    \end{itemize}
    \column{0.6\textwidth}
    \figlivroslide{fig_21_27.jpg}{0.5\textheight}{Fig.~21.27 --- os
      índices do ENSO oscilando pelas décadas}
  \end{columns}
\end{frame}

% ---------------------------------------------------------------
\section{Síntese e laboratório}

\begin{frame}{O mapa do capítulo}
  \eqdestaque{$\Delta T = \tfrac{\lambda_0\Delta F}{1-f}$ \quad
    gelo-albedo: 2 equilíbrios \quad $Q_{65N}$(órbita) \quad
    $\dot T = aT + bh$}
  \vspace{0.6em}
  Feedbacks $+$ escalas lentas $+$ forçantes externas $=$ um clima que
  varia sozinho — a régua contra a qual o Cap.~22 mede a perturbação
  humana.
\end{frame}

\begin{frame}{Laboratório numérico --- notebook do Cap.~21}
  \texttt{notebooks/cap21\_clima.ipynb}
  \begin{enumerate}
    \item Bifurcação gelo-albedo (bola de neve).
    \item Milankovitch e a insolação de 65°N.
    \item Daisyworld completo.
    \item O oscilador de recarga do ENSO.
  \end{enumerate}
  \vspace{0.4em}
  \textbf{Exercícios}: T1--T5 e N1--N4 nas notas; células reservadas no
  notebook. Os quatro casos são papers clássicos em miniatura —
  Budyko, Milankovitch, Watson--Lovelock, Jin.
  \vfill
  \atribuicao
\end{frame}

\end{document}
